\chapter{Kommentierte Quellen}

Achtundzwanzig Werke, eingeordnet: Zu jedem steht, wofür man es aufschlägt und was es leistet, das andere nicht leisten.

\lit{1}{Angrist, J., Pischke, J.-S.: \textit{Mostly Harmless Econometrics}. Princeton 2009.}{Die Verfahren der Kapitel~9 und 10 aus ökonometrischer Sicht --- Instrumente, Sprünge, doppelte Differenzen. Salopp geschrieben, technisch genau, mit vielen Beispielen aus der Arbeitsmarktforschung. Ergänzt diesen Band dort, wo er am knappsten ist.}

\lit{2}{Athey, S., Wager, S.: \textit{Estimating Treatment Effects with Causal Forests}. Observational Studies 7, 2021.}{Kausale Wälder aus Kapitel~14, von ihren Entwicklern erklärt. Zeigt, wie Bäume nach Effektunterschieden statt nach Ausgängen teilen und warum die Trennung von Teilen und Schätzen nötig ist.}

\lit{3}{Baron, R., Kenny, D.: \textit{The moderator--mediator variable distinction}. J.\ Personality Soc.\ Psych.\ 51, 1986.}{Der meistzitierte Aufsatz zur Mediation --- und ein Lehrstück darüber, wie lange ein bequemes Verfahren überlebt. \textbf{Zu lesen, um zu verstehen, was in Kapitel~13 berichtigt wird}, nicht als Anleitung.}

\lit{4}{Bradford Hill, A.: \textit{The Environment and Disease: Association or Causation?} Proc.\ R.\ Soc.\ Med.\ 58, 1965.}{\textbf{Die neun Gesichtspunkte} zur Beurteilung, ob ein Zusammenhang ursächlich ist: Stärke, Beständigkeit, Dosis-Wirkung und sechs weitere. Hill nannte sie ausdrücklich keine Prüfliste. Sie werden bis heute als solche benutzt und haben die Rauchen-Krebs-Debatte entschieden.}

\lit{5}{Cinelli, C., Forney, A., Pearl, J.: \textit{A Crash Course in Good and Bad Controls}. Sociol.\ Methods Res.\ 53, 2024.}{Kapitel~5 dieses Bandes in Aufsatzform: achtzehn Graphen, jeder mit der Antwort, ob man anpassen darf. \textbf{Zum Aushängen über dem Schreibtisch.} Enthält auch die M-Verzerrung und die Verstärkung durch Instrumente.}

\lit{6}{Cornfield, J. u.\,a.: \textit{Smoking and lung cancer}. J.\ Natl.\ Cancer Inst.\ 22, 1959.}{Die Antwort auf Fishers Einwand aus Kapitel~1, und zugleich die Geburt der Empfindlichkeitsanalyse: Eine verborgene Ursache hätte das Risiko verneunfachen müssen. Ein Muster dafür, wie man ohne Randomisierung überzeugend argumentiert.}

\lit{7}{Cunningham, S.: \textit{Causal Inference: The Mixtape}. Yale 2021.}{Eine ungewöhnlich lesbare Einführung, mit Code und historischen Abschweifungen. Deckt die Kapitel~7 bis 10 vollständig ab. Frei verfügbar; gut für Leser, die selbst rechnen wollen.}

\lit{8}{Elwert, F., Winship, C.: \textit{Endogenous Selection Bias}. Annu.\ Rev.\ Sociol.\ 40, 2014.}{Der Kollider, systematisch --- als Auswahl, als Anpassung, als Studieneinschluss. Erklärt, warum dieselbe Verzerrung unter einem Dutzend Namen umläuft, und ist die beste Vertiefung zu Aufgabe 1.2 und 5.2.}

\lit{9}{Greenland, S., Pearl, J., Robins, J.: \textit{Causal Diagrams for Epidemiologic Research}. Epidemiology 10, 1999.}{Der Aufsatz, der die Graphen in die Epidemiologie trug --- zehn Seiten, ohne Beweise, mit Beispielen. \textbf{Wer nur einen Text über kausale Graphen liest, lese diesen.} Enthält die M-Verzerrung, die die alte Anpassungsregel widerlegte.}

\lit{10}{Hernán, M., Robins, J.: \textit{Causal Inference: What If}. Chapman \& Hall 2020.}{\textbf{Das Standardwerk und die Hauptquelle dieses Bandes.} Teil I ohne Modelle, Teil II mit, Teil III über die Zeit --- g-Formel, Gewichtung, Zielstudie. Von den Entwicklern der Verfahren, medizinisch orientiert, frei verfügbar.}

\lit{11}{Hernán, M. u.\,a.: \textit{Observational studies analyzed like randomized experiments}. Epidemiology 19, 2008.}{Die Auflösung des Hormonwiderspruchs aus Kapitel~11. Zeigt an einem Fall von großer Tragweite, dass die meisten Widersprüche zwischen Beobachtung und Versuch aus dem Entwurf stammen, nicht aus Störgrößen.}

\lit{12}{Hernán, M., Robins, J.: \textit{Using Big Data to Emulate a Target Trial}. Am.\ J.\ Epidemiol.\ 183, 2016.}{Die sieben Zeilen aus Kapitel~11, ausgeführt. Kurz und unmittelbar anwendbar. \textbf{Vor jeder Auswertung von Routinedaten zu lesen} --- und auszufüllen.}

\lit{13}{Holland, P.: \textit{Statistics and Causal Inference}. J.\ Am.\ Stat.\ Assoc.\ 81, 1986.}{Der Aufsatz, der das grundlegende Problem benannte und die Diskussion um \enquote{keine Ursache ohne Manipulation} auslöste. Mit Erwiderungen, die den Streit gut abbilden. Philosophisch klar, mathematisch leicht.}

\lit{14}{Imbens, G., Angrist, J.: \textit{Identification and Estimation of Local Average Treatment Effects}. Econometrica 62, 1994.}{Satz 9.2 in seiner Urform: Der Instrumentschätzer gilt für die Befolger. \textbf{Eine unwillkommene Nachricht, die ein ganzes Fach zum Umdenken zwang} --- und 2021 mit dem Nobelpreis bedacht wurde.}

\lit{15}{Imbens, G., Rubin, D.: \textit{Causal Inference for Statistics, Social, and Biomedical Sciences}. Cambridge 2015.}{Das große Werk zur Sprache der potentiellen Ergebnisse --- Randomisierung, Neigung, Instrumente, mit vollständigen Herleitungen. Ohne Graphen; die Ergänzung zu [18] von der anderen Schule.}

\lit{16}{Kuß, O., Blettner, M., Börgermann, J.: \textit{Propensity Score --- eine alternative Methode zur Analyse von Therapieeffekten}. Dtsch.\ Ärztebl.\ 113, 2016.}{\textbf{Auf Deutsch und kurz.} Die Behandlungsneigung aus Kapitel~7, für Ärzte erklärt, mit einem herzchirurgischen Beispiel. Frei verfügbar; gut als erste Lektüre oder zur Weitergabe an Kollegen.}

\lit{17}{Mill, J.\,S.: \textit{A System of Logic}. London 1843.}{Die Methoden der Übereinstimmung und des Unterschieds --- der erste systematische Versuch, Ursachen aus Beobachtung zu erschließen. Mills \enquote{Methode des Unterschieds} ist im Kern die randomisierte Studie, hundert Jahre zu früh gedacht.}

\lit{18}{Pearl, J.: \textit{Causality: Models, Reasoning, and Inference}. 2.\ Aufl., Cambridge 2009.}{\textbf{Das Werk, das die Graphen begründete.} Enthält d-Trennung, Hintertür, Vordertür, do-Kalkül mit Beweisen --- die Fundstelle für alle Auslassungen der Kapitel~4 und 6. Anspruchsvoll und eigenwillig geschrieben.}

\lit{19}{Pearl, J., Glymour, M., Jewell, N.: \textit{Causal Inference in Statistics: A Primer}. Wiley 2016.}{Die zugängliche Fassung von [18]: knapp zweihundert Seiten, mit Aufgaben und Lösungen. \textbf{Der beste Einstieg in die grafische Seite des Faches} --- für Leser, die Kapitel~4 bis 6 vertiefen wollen.}

\lit{20}{Robins, J.: \textit{A new approach to causal inference in mortality studies}. Math.\ Modelling 7, 1986.}{Die Arbeit über die Kupferhütte, in der die g-Formel entstand. $115$ Seiten, schwer zugänglich --- aber sie enthält alles, was Kapitel~12 behandelt, zwanzig Jahre vor seiner allgemeinen Anerkennung.}

\lit{21}{Rosenbaum, P., Rubin, D.: \textit{The central role of the propensity score}. Biometrika 70, 1983.}{Satz 7.2 im Original --- acht Seiten, die die Praxis der Beobachtungsforschung verändert haben. Der Beweis ist so kurz wie der in Kapitel~7.}

\lit{22}{Rubin, D.: \textit{Estimating causal effects of treatments in randomized and nonrandomized studies}. J.\ Educ.\ Psychol.\ 66, 1974.}{Die Übertragung der potentiellen Ergebnisse von Versuchen auf Beobachtungsdaten. Der Ausgangspunkt der gesamten Rubin-Schule; bemerkenswert klar geschrieben.}

\lit{23}{Schomaker, M., Heumann, C.: \textit{Kausalität in der Statistik}. Springer Spektrum 2023.}{\textbf{Auf Deutsch und mathematisch.} Behandelt Graphen, potentielle Ergebnisse und die Schätzverfahren in einem Band, mit Beispielen und R-Code. Die deutschsprachige Entsprechung zu [10] in kompakterer Form.}

\lit{24}{Shpitser, I., Pearl, J.: \textit{Identification of joint interventional distributions}. AAAI 2006.}{Der Vollständigkeitsbeweis des do-Kalküls: Lässt sich ein Effekt bestimmen, so führen die drei Regeln zum Ziel. Technisch; wichtig als Beleg dafür, dass \enquote{unbestimmbar} eine beweisbare Aussage ist.}

\lit{25}{Suttorp, M. u.\,a.: \textit{Graphical presentation of confounding in directed acyclic graphs}. Nephrol.\ Dial.\ Transplant.\ 30, 2015.}{Kausale Graphen für klinisch tätige Ärzte, mit Beispielen aus der Nephrologie. Kurz, praktisch, ohne Formalismus --- gut geeignet, um Kollegen den Unterschied zwischen Störgröße, Mediator und Kollider zu erklären.}

\lit{26}{VanderWeele, T.: \textit{Explanation in Causal Inference: Methods for Mediation and Interaction}. Oxford 2015.}{Das Standardwerk zu Kapitel~13 --- natürliche und kontrollierte Effekte, Wechselwirkungen, Empfindlichkeitsanalysen. Die Fundstelle für den Satz über natürliche Effekte, den dieser Band ausließ.}

\lit{27}{VanderWeele, T., Ding, P.: \textit{Sensitivity Analysis in Observational Research: Introducing the E-Value}. Ann.\ Intern.\ Med.\ 167, 2017.}{Eine einzige Zahl, die angibt, wie stark eine unbeobachtete Störgröße wirken müsste, um den Befund zu erklären. \textbf{Sie gehört in jede Beobachtungsstudie} und ist ohne Annahmen über die Form der Störung berechenbar.}

\lit{28}{Wong, A. u.\,a.: \textit{External validation of a widely implemented proprietary sepsis prediction model}. JAMA Intern.\ Med.\ 181, 2021.}{Die Nachprüfung aus Band 22, hier als Mahnung am anderen Ende: Ein Modell, das vorhersagt, ohne Ursachen zu kennen, bricht zusammen, sobald sich die Verteilung ändert. Der Übergang von jenem Band zu diesem in einer Arbeit.}

\begin{merksatz}
\boxtitel{Wenn Sie nur drei Werke lesen}
\begin{center}
\renewcommand{\arraystretch}{1.6}
\begin{tabular}{@{}P{1.0cm}P{3.4cm}P{3.8cm}@{}}
\toprule
\textbf{[10]} & \textbf{Hernán, Robins} & \textbf{Für alles}\\
\textbf{[9]} & \textbf{Greenland u.\,a.} & Für die Graphen\\
\textbf{[12]} & \textbf{Hernán, Robins} & \textbf{Vor der ersten Auswertung}\\
\bottomrule
\end{tabular}
\end{center}

\begin{center}
\begin{minipage}{0.88\textwidth}
\kursiv{Hernán und Robins erklären die Verfahren; Greenland und Mitarbeiter liefern auf zehn Seiten, wofür Kapitel~4 bis 6 fünfzig brauchten. \textbf{Und der dritte Text ist die Prüfliste, die die meisten Fehler verhindert, bevor sie entstehen.} Alle drei sind frei verfügbar.}
\end{minipage}
\end{center}

\begin{center}
\begin{minipage}{0.88\textwidth}
\kursiv{Nach Interessenlage: [19] Pearl u.\,a.\ zum Einstieg in die Graphen, [1] Angrist und Pischke für natürliche Versuche, [26] VanderWeele für Mechanismen, [27] für die Empfindlichkeit. Auf Deutsch: [23] Schomaker und Heumann für die Mathematik, [16] Kuß u.\,a.\ und [25] Suttorp u.\,a.\ für die klinische Praxis.}
\end{minipage}
\end{center}
\end{merksatz}

\bigskip
\begin{flushright}
\itshape\color{grau}
Die Daten sagen, was zusammen auftritt.\\
Der Graph sagt, was man annimmt.\\
Erst beides zusammen sagt,\\
was geschieht, wenn man handelt.
\end{flushright}
